% homography_explained.tex
% Standalone, self-contained explainer. Compile with:
%     pdflatex homography_explained.tex
% (no bibliography, no external figures required)
\documentclass[11pt]{article}

\usepackage[margin=1in]{geometry}
\usepackage{amsmath,amssymb}
\usepackage{booktabs}
\usepackage{array}
\usepackage{xcolor}
\usepackage{listings}
\usepackage[colorlinks=true,linkcolor=blue,urlcolor=blue]{hyperref}

% "vex" operator (skew-symmetric matrix -> 3-vector)
\DeclareMathOperator{\vex}{vex}

% compact inline matrices for use inside table cells
\newcommand{\sm}[1]{\left[\begin{smallmatrix}#1\end{smallmatrix}\right]}

% C++ code listing style
\definecolor{codegray}{gray}{0.95}
\definecolor{codekw}{rgb}{0.0,0.0,0.6}
\definecolor{codecom}{rgb}{0.0,0.5,0.0}
\lstset{
  language=C++,
  basicstyle=\ttfamily\small,
  keywordstyle=\color{codekw}\bfseries,
  commentstyle=\color{codecom}\itshape,
  backgroundcolor=\color{codegray},
  frame=single,
  framerule=0pt,
  breaklines=true,
  showstringspaces=false,
  columns=fullflexible,
  xleftmargin=6pt,
}

\title{Homography, Affine, Planar --- Explained From Scratch\\[2pt]
\large for the \texttt{h\_vs} controller}
\author{}
\date{}

\begin{document}
\maketitle

\noindent\textbf{Who this is for:} you've taken linear algebra --- vectors,
matrices, matrix multiplication, inverse, transpose, the identity matrix.
That's all you need. No projective geometry, no Lie groups, no differential
geometry. We build everything else from those pieces and then read the actual
\texttt{h\_vs} code with it.

% ─────────────────────────────────────────────────────────────────────────────
\section{The one-sentence idea}

\begin{quote}
A \textbf{homography} is a single $3\times3$ matrix that describes how a
\textbf{flat surface} (a plane --- e.g.\ the face of a stop sign) appears to
move/warp between two camera views.
\end{quote}

If you know where the four corners of the sign are \emph{now}, and where you
\emph{want} them to be, the homography is the matrix that connects those two
pictures. The \texttt{h\_vs} controller measures that matrix and turns ``how
far is the matrix from doing nothing'' into ``how should the drone move.''

That's the whole story. The rest is unpacking three words: \textbf{homogeneous},
\textbf{affine}, \textbf{planar}.

% ─────────────────────────────────────────────────────────────────────────────
\section{The ``add a 1'' trick: homogeneous coordinates}

A pixel is a 2D point $(u,v)$. We staple a $1$ on the end and treat it as a 3D
vector:
\[
(u,v)\;\longrightarrow\;\begin{bmatrix} u\\ v\\ 1\end{bmatrix}.
\]
Why? Two reasons, both pure linear algebra.

\paragraph{Reason A --- translation becomes multiplication.}
Normally ``shift right by $t_x$, down by $t_y$'' is an \emph{addition}. With the
extra $1$ it becomes a matrix multiply:
\[
\begin{bmatrix} 1&0&t_x\\ 0&1&t_y\\ 0&0&1\end{bmatrix}
\begin{bmatrix} u\\ v\\ 1\end{bmatrix}
=\begin{bmatrix} u+t_x\\ v+t_y\\ 1\end{bmatrix}.
\]
Now \emph{every} operation --- rotate, scale, shift, even perspective --- is one
$3\times3$ matrix. That uniformity is the entire point.

\paragraph{Reason B --- we may rescale the whole vector.}
In this homogeneous world these all mean the \emph{same} pixel:
\[
\begin{bmatrix} u\\ v\\ 1\end{bmatrix}\sim
\begin{bmatrix} 2u\\ 2v\\ 2\end{bmatrix}\sim
\begin{bmatrix} 10u\\ 10v\\ 10\end{bmatrix}.
\]
The rule to recover a real pixel: \textbf{divide by the last entry.} So
$(2u,2v,2)\to(u,v)$. This ``divide by the last number'' is exactly what produces
\emph{perspective} (far things look smaller). Because of this freedom a
homography only matters \textbf{up to scale}: multiplying the whole matrix by
$7$ gives the same transform (hence people normalize the bottom-right entry
to~$1$).

% ─────────────────────────────────────────────────────────────────────────────
\section{The ladder of transformations (rigid $\to$ projective)}

All of these are $3\times3$ matrices acting on $(u,v,1)$. They differ only in
how much freedom the matrix has.

\begin{center}
\renewcommand{\arraystretch}{1.6}
\begin{tabular}{@{}lccl@{}}
\toprule
Name & Matrix shape & Can do & Preserves \\
\midrule
Translation & $\sm{1&0&t_x\\0&1&t_y\\0&0&1}$ & slide & all but position \\
Euclidean (rigid) & $\sm{R&\;t\\0\,0&1}$ & rotate $+$ slide & lengths, angles \\
Similarity & $\sm{sR&\;t\\0\,0&1}$ & $+$ uniform zoom & angles, ratios \\
Affine & $\sm{a&b&t_x\\c&d&t_y\\0&0&1}$ & $+$ shear, non-uniform scale & \textbf{parallel lines} \\
Projective (homography) & $\sm{a&b&t_x\\c&d&t_y\\ \mathbf{g}&\mathbf{h}&1}$ & $+$ perspective & \textbf{straight lines only} \\
\bottomrule
\end{tabular}
\end{center}

The single most important row is the \textbf{bottom row}.

\begin{itemize}
\item \textbf{Affine} has bottom row $[0,\,0,\,1]$. The last entry of the output
stays exactly $1$, so ``divide by the last entry'' does nothing.
\textbf{No perspective.} Parallel lines in, parallel lines out. Think of a
photocopier: a flat page scanned straight-on, scaled/sheared but never tilted
away.
\item \textbf{Projective / homography} has a general bottom row $[g,h,1]$. The
last entry of the output is $g\,u+h\,v+1$, which \textbf{changes from pixel to
pixel}. Dividing by it squishes some points toward the center --- that's
perspective foreshortening (railroad tracks converging to a vanishing point).
\end{itemize}

\begin{quote}
\textbf{So ``affine'' is just ``a homography that forgot how to do
perspective.''} It is the special case $g=h=0$. Keep this sentence; it is the
key to the whole \texttt{h\_vs} story in Section~\ref{sec:affinetrap}.
\end{quote}

\paragraph{Tiny numeric example of the perspective divide.}
Take the projective matrix with $g=0.001$, $h=0$:
\[
\begin{bmatrix} 1&0&0\\ 0&1&0\\ 0.001&0&1\end{bmatrix}
\begin{bmatrix} 500\\ 200\\ 1\end{bmatrix}
=\begin{bmatrix} 500\\ 200\\ 1.5\end{bmatrix}
\xrightarrow{\;\div\,1.5\;}
\begin{bmatrix} 333\\ 133\end{bmatrix}.
\]
The point got pulled inward, and \emph{how much} depended on its
$u$-coordinate~($500$). An affine matrix ($g=h=0$) could never do that --- it
leaves the last entry at $1$ and pulls nothing.

% ─────────────────────────────────────────────────────────────────────────────
\section{What ``planar'' means and why it's required}

A homography is only \emph{exactly} correct for points that all lie on
\textbf{one flat plane} in the real world. A stop-sign face is a plane. A wall
is a plane. The ground is a plane.

The theorem (proof not needed): \textbf{two images of the same flat plane are
always related by a single $3\times3$ homography.} Four corner-pairs is exactly
enough to solve for it --- each corner gives $2$ equations, $H$ has $8$ free
numbers, and $4\times2=8$. That's why the controller uses the \textbf{four
bounding-box corners} of the detected sign: a box is planar, and $4$ corners
pin down $H$ exactly.

% ─────────────────────────────────────────────────────────────────────────────
\section{The camera matrix $K$ (pixels $\leftrightarrow$ rays)}

A camera turns a 3D direction into a pixel via the \emph{intrinsics} matrix
\[
K=\begin{bmatrix} f_x&0&c_{x0}\\ 0&f_y&c_{y0}\\ 0&0&1\end{bmatrix},
\]
where $f_x,f_y$ are focal lengths in pixels (how ``zoomed in'' the camera is)
and $(c_{x0},c_{y0})$ is the \emph{principal point} (the pixel at the optical
center, essentially the middle of the image).

$K$ maps a normalized 3D ray $\to$ a pixel; so $K^{-1}$ maps a
\textbf{pixel $\to$ a normalized ray} (a real-world direction, stripped of the
camera's zoom/offset). Remember: \textbf{$K^{-1}$ converts ``where it is on
screen'' into ``which real-world direction it is.''} We need that to talk about
drone motion in meters/radians instead of pixels.

% ─────────────────────────────────────────────────────────────────────────────
\section{The control law, exactly as the code computes it}
\label{sec:law}

Here is the whole math core from
\texttt{src/h\_vs\_servo/src/homography\_2d\_vs.cpp}:

\begin{lstlisting}
Eigen::Matrix3d H = _K.inverse() * G * _K;          // line 15
Eigen::Vector3d m_star = _K.inverse() * p_star;     // line 16
twist << _lambda_v.asDiagonal() * _computeEv(H, m_star),   // (H - I) m*
         _lambda_w.asDiagonal() * _computeEw(H);           // vex(H - H^T)
\end{lstlisting}

\subsection{Measure $G$ (pixel-space homography)}
In \texttt{hvs\_controller.cpp} the controller builds two sets of four corners:
\texttt{pts\_star\_} (where we \emph{want} the box --- centered, target size)
and \texttt{pts\_curr} (where the box \emph{is}). OpenCV solves the 4-corner
system:
\begin{lstlisting}
cv::Mat G_cv = cv::getPerspectiveTransform(pts_star_, pts_curr);  // line 72
\end{lstlisting}
$G$ is the $3\times3$ mapping \emph{desired pixels $\to$ current pixels}.

\subsection{Convert to the ``Euclidean'' homography $H$}
\[
\boxed{\,H = K^{-1}\,G\,K\,}
\]
This is a \emph{change of basis} (sandwich a matrix between $K^{-1}$ and $K$).
It re-expresses the same warp in normalized-ray coordinates instead of pixels,
so $H$ now has physical meaning.

\begin{quote}
\textbf{Key intuition:} if the drone is \emph{exactly} where it should be, then
current $=$ desired, $G$ does nothing, and $H=I$ (the identity). \textbf{The
error is ``how far is $H$ from the identity.''} The next two formulas are just
two different ways to measure ``$H$ minus doing-nothing.''
\end{quote}

\subsection{Translation error --- how far off in \emph{position}}
\[
\mathbf{e}_v=(H-I)\,\mathbf{m}^{*}\qquad\text{(paper's eq.\ 15)}
\]
The code defaults $\mathbf{m}^{*}$ to the principal point,
$p^{*}=(c_{x0},c_{y0},1)$, so $\mathbf{m}^{*}=K^{-1}p^{*}=(0,0,1)$ --- ``straight
ahead.'' Multiplying $(H-I)$ by $(0,0,1)$ simply \textbf{picks out the third
column of $H-I$}:
\[
\mathbf{e}_v=\begin{bmatrix} H_{0,2}\\ H_{1,2}\\ H_{2,2}-1\end{bmatrix}
=\begin{bmatrix}\text{sideways offset}\\ \text{up/down offset}\\ \text{depth
offset}\end{bmatrix}.
\]

\subsection{Rotation error --- how far off in \emph{orientation}}
\[
\mathbf{e}_\omega=\vex\!\left(H-H^{\top}\right)\qquad\text{(paper's eq.\ 16)}
\]
One fact you can prove with transpose rules: for \textbf{any} matrix $H$, the
matrix $S=H-H^{\top}$ is \textbf{skew-symmetric}, i.e.\ $S^{\top}=-S$:
\[
(H-H^{\top})^{\top}=H^{\top}-H=-(H-H^{\top}).\;\checkmark
\]
A skew-symmetric $3\times3$ always has zeros on the diagonal and only
\textbf{three} independent numbers:
\[
S=\begin{bmatrix} 0&-z&y\\ z&0&-x\\ -y&x&0\end{bmatrix}.
\]
The operator $\vex(\cdot)$ (``vee'') just \textbf{reads those three numbers out}
into the vector $(x,y,z)$ --- the inverse of building the skew matrix. In code:
\begin{lstlisting}
return Eigen::Vector3d(H_skew(2,1), H_skew(0,2), H_skew(1,0));   // line 39
\end{lstlisting}
Those three numbers are the small camera rotation (roll/pitch/yaw) needed to
line back up with the sign. \emph{Why does the antisymmetric part encode
rotation?} Because a tiny rotation matrix is $I+S$ with $S$ skew-symmetric;
subtracting $H^{\top}$ cancels the symmetric stretch part and leaves the
rotation.

\subsection{Stack into a 6-DOF twist, with gains}
\[
\text{twist}=\begin{bmatrix}\boldsymbol{\lambda}_v\odot\mathbf{e}_v\\
\boldsymbol{\lambda}_\omega\odot\mathbf{e}_\omega\end{bmatrix}
=[\,t_x,t_y,t_z,\;r_x,r_y,r_z\,],
\]
where $\odot$ is element-wise multiply (the gains scale each axis). The gains
live in \texttt{config/hil/bench\_h\_vs.yaml}:
$\boldsymbol{\lambda}_v=[1.0,1.0,0.0]$,
$\boldsymbol{\lambda}_\omega=[0.0,0.4,0.0]$.

% ─────────────────────────────────────────────────────────────────────────────
\section{From camera twist to drone motion (the body mapping)}

\texttt{hvs\_controller.cpp} (lines 112--117) converts the camera twist into
body velocities, after a \textbf{5-tick moving average} (averaging the last $5$
computations kills jitter --- this is exactly why \texttt{h\_vs} has the
smoothest commands in the benchmark):

\begin{lstlisting}
v.vx = std::max(0.0, k_fwd_ * (target_bbox_ratio_ - in.bbox_ratio)); // forward
v.vy = -avg(0);   // camera "off to the side" -> body left/right
v.vz = -avg(1);   // camera "up/down"          -> body up/down
v.wz = -avg(4);   // camera yaw (e_w[1])       -> body yaw
\end{lstlisting}

Note that \texttt{vx} does \textbf{not} come from the homography --- it is a plain
proportional law on bounding-box size (\texttt{bbox\_ratio}, the box height
fraction, a stand-in for range). The next section explains \emph{why it has to
be that way here}.

% ─────────────────────────────────────────────────────────────────────────────
\section{Why the simulator's box breaks the homography (the ``affine trap'')}
\label{sec:affinetrap}

This is the crux, and the issue flagged in the paper and report.

In the benchmark, detection comes from an \textbf{oracle} that draws a
\textbf{perfectly axis-aligned} rectangle (sides parallel to the image edges,
never tilted). Mapping four axis-aligned corners to four axis-aligned corners can
be satisfied \textbf{without any perspective} --- the solution is
\textbf{affine}. From Section~3, affine means bottom row $[0,0,1]$, so after
$H=K^{-1}GK$:
\[
H_{2,2}=1\quad\text{always.}
\]
Two consequences fall straight out of the formulas in Section~\ref{sec:law}.

\paragraph{(i) No forward motion.}
The depth component of the translation error is
$\mathbf{e}_v[2]=H_{2,2}-1=0$, \emph{always}, no matter how far the drone is.
The homography simply cannot sense or command range here, so the code bypasses
it with the proportional $v_x=k_{\mathrm{fwd}}(\rho^{*}-\rho)$ (line 113). With a
real, textured, tilted target the corners would \emph{not} be axis-aligned, $G$
would be genuinely projective, $H_{2,2}\neq1$, and the homography would produce
forward motion on its own.

\paragraph{(ii) ``Lateral parking.''}
For an affine $H$, the two horizontal error terms are
\[
\underbrace{\mathbf{e}_v[0]=H_{0,2}=\frac{c_x-c_{x0}}{f_x}}_{\text{slide
sideways }(v_y)}
\qquad
\underbrace{\mathbf{e}_\omega[1]=H_{0,2}-H_{2,0}=\frac{c_x-c_{x0}}{f_x}-0}_{\text{yaw
}(w_z)}.
\]
They are the \textbf{same number} --- the horizontal pixel offset of the sign.
So the \emph{same} error drives both ``slide sideways'' and ``rotate.'' Rotating
is the cheaper way to make the sign appear centered, so the drone yaws to face
the sign from wherever it is --- and the instant the sign looks centered,
\emph{both} corrections switch off together. Result: it parks slightly
\textbf{to the side} of the sign, facing it diagonally, instead of squaring up
in front.

This is a degeneracy of the \emph{input} (a flat, untextured, axis-aligned box
carries no perspective information to tell ``I'm off to the side'' apart from
``I'm rotated''), \textbf{not a bug in the code}.

% ─────────────────────────────────────────────────────────────────────────────
\section{The paper's symbols, decoded}

\begin{center}
\renewcommand{\arraystretch}{1.4}
\begin{tabular}{@{}lll@{}}
\toprule
Symbol & Reads as & Plain meaning \\
\midrule
$\mathbf{G}$ & ``gee'' & pixel-space homography (desired $\to$ current corners) \\
$\mathbf{K}$ & ``kay'' & camera intrinsics (focal lengths $+$ principal point) \\
$H=K^{-1}GK$ & & same warp in real-world ray coords; $=I$ when on-target \\
$\mathbf{m}^{*}$ & ``em star'' & desired reference ray $K^{-1}p^{*}=(0,0,1)$, straight ahead \\
$\mathbf{e}_v=(H-I)\mathbf{m}^{*}$ & & translation error $=$ 3rd column of $H-I$ \\
$\mathbf{e}_\omega=\vex(H-H^{\top})$ & & rotation error $=$ the 3 numbers of the skew part \\
$H_{2,2}=1\Rightarrow e_{v,z}=0$ & & affine box $\Rightarrow$ no depth error $\Rightarrow$ no forward push \\
$(c_x-c_{x0})/f_x$ & & horizontal pixel offset (drives both $v_y$ and $w_z$) \\
\bottomrule
\end{tabular}
\end{center}

When the paper calls the homography law \emph{image-based} yet says it
\emph{``places it between pure image-based and pose-based control''}: it only
ever measures things in the image (the four corners), but because $H=K^{-1}GK$
secretly contains rotation $+$ scaled-translation structure, it borrows a little
of the ``reconstruct the pose'' flavor --- hence ``between.''

% ─────────────────────────────────────────────────────────────────────────────
\section{File map}

\begin{center}
\renewcommand{\arraystretch}{1.3}
\begin{tabular}{@{}p{0.45\textwidth}p{0.40\textwidth}c@{}}
\toprule
Concept & File & Lines \\
\midrule
Solve $G$ from 4 corners; build desired corners & \texttt{h\_vs\_servo/src/hvs\_controller.cpp} & 43--72 \\
$H=K^{-1}GK$, $\mathbf{e}_v$, $\mathbf{e}_\omega$, twist & \texttt{h\_vs\_servo/src/homography\_2d\_vs.cpp} & 12--40 \\
$\vex$ (read 3 numbers from skew matrix) & \texttt{homography\_2d\_vs.cpp} & 37--40 \\
Twist $\to$ body velocities; 5-tick average & \texttt{hvs\_controller.cpp} & 87--117 \\
Gains $\boldsymbol{\lambda}_v,\boldsymbol{\lambda}_\omega$, \texttt{k\_fwd}, buffer & \texttt{config/hil/bench\_h\_vs.yaml} & 1--19 \\
Paper write-up $+$ affine-degeneracy caveats & \texttt{paper.tex} & \texttt{sec:vs\_arch}, \texttt{sec:vs\_discuss} \\
\bottomrule
\end{tabular}
\end{center}

% ─────────────────────────────────────────────────────────────────────────────
\section{One-paragraph recap}

Staple a $1$ onto pixels so every transform is a $3\times3$ matrix. A
\textbf{homography} is the matrix relating two views of a \textbf{flat plane};
an \textbf{affine} transform is the crippled homography whose bottom row is
$[0,0,1]$, so it cannot do perspective. \texttt{h\_vs} measures the homography
$G$ from the four box corners, rewrites it as $H=K^{-1}GK$ (which equals the
identity when the drone is on-target), and reads the \emph{position} error off
the third column of $H-I$ and the \emph{orientation} error off the
skew-symmetric part $H-H^{\top}$. Because the simulator's boxes are
axis-aligned, the measured homography collapses to affine, which zeroes the
depth term (no forward push from the homography $\to$ handled by a separate
proportional law) and makes the sideways and yaw errors identical ($\to$ the
controller parks beside the sign). With a real, textured, tilted target the full
perspective returns and those degeneracies disappear.

\end{document}
